%!TEX program = xelatex
%!TEX root = ../all.tex
\chapter{Latent variable models}

\section*{Inference and Maximum Likelihood Recall}\label{sec:ml}

The construction of a probabilistic model from a finite sample of observed data is one of
the central and most fundamental problems in pattern recognition and machine learning.
Given a collection of observations, we wish to identify a probability distribution that
faithfully describes the process that generated them---one that not only fits the data at
hand, but generalises to unseen instances drawn from the same underlying process.

A standard and widely adopted approach to this problem involves \emph{parametric
models}: rather than searching over all conceivable probability distributions, one proposes
a specific functional form for the density, indexed by a finite-dimensional vector of
parameters. The problem of model construction then reduces to the problem of estimating
these parameters from data. This reduction is both a strength and a limitation: it makes
the problem tractable, but it introduces a modelling assumption whose validity must be
justified by domain knowledge or empirical evaluation.

Concretely, we assume that the observed data $\ox$ is a realisation of a random variable
$\x$ defined on a domain $\mathcal{X}$, distributed according to a parametric family
$p(\x; \Vtheta)$ with parameter vector $\Vtheta = (\theta_1, \ldots, \theta_m)$.\footnote{We
may interpret $p(\x; \Vtheta)$ as a function of $\x$ given the value of the parameter
$\Vtheta$, or as a function of $\Vtheta$ given an observed value of $\x$.} The true
distribution of $\x$ is unknown; we only commit to the assumption that it belongs to the
parametric family $\{p(\x;\vrt) : \vrt \in \Theta\}$ for some parameter space $\Theta$.

This gives rise to an \alert{inference} problem: given the observation $\ox$, find the
parameter value $\prts$ such that $p(\x;\prts)$ best represents the data. The qualifier
``best'' requires a criterion, and the most natural and widely used one is to choose the
parameter value that assigns the highest probability to the observed data. This approach is
called \alert{point estimation}, as it produces a single estimate $\prts$ rather than a full
posterior distribution over parameters.

\subsection*{Maximum Likelihood Estimation}

Given a parameter value $\prt$, the \alert{evidence} $p(\vlx;\prt)$ of the observation
$\vlx$ is simply the probability that the model assigns to $\vlx$ under parameter $\prt$.
When viewed as a function of $\Vtheta$ for a fixed observation $\vlx$, this quantity is
called the \alert{likelihood} $L(\vrt|\vlx)$. The likelihood has a natural and intuitive
interpretation: it measures how plausible the parameter value $\vrt$ is, given what we have
observed. A high likelihood means the model with parameter $\vrt$ would have been very
likely to produce the observed data; a low likelihood indicates the opposite.

The \alert{maximum likelihood estimate} (MLE) is defined as the parameter value that
maximises the likelihood:
\begin{myequation}
	\prts=\argmax{\Vtheta} L(\vrt|\vlx).
\end{myequation}
Since the logarithm is a strictly monotone increasing function, maximising the likelihood is
equivalent to maximising its logarithm:
\begin{myequation}
	\val\prts=\argmax{\Vtheta} \log L(\Vtheta|\ox)=\argmax{\Vtheta} l(\vrt|\vlx),
\end{myequation}
where $l(\vrt|\vlx)\defeq\log L(\vrt|\vlx)$ is the \alert{log-likelihood} of
$\vlx$.\footnote{The equality holds because both functions have the same stationary
points, since $\frac{d}{dx}\log f(x)=\frac{1}{f(x)}\frac{d}{dx}f(x)$ and $f(x)$ is
bounded, and because $f(x_1)>f(x_2)$ if and only if $\log f(x_1)>\log f(x_2)$.}
Working with the log-likelihood is almost always preferable in practice: it transforms
products into sums (which are numerically more stable), and it simplifies the computation
of gradients and stationarity conditions.

Finding the MLE proceeds differently depending on the complexity of the model. In the
simplest and most favourable case, the stationarity condition\footnote{This corresponds to
the system of equations $\pdv{l(\vrt|\vlx)}{\vr{\theta_i}}=0$ for $i=1,\ldots,n$.}
\begin{myequation}
	\gradient{\vrt} {l(\vrt|\vlx)}=\gradient{\vrt}{\log p(\vlx|\vrt)}=\Zero
\end{myequation}
admits a finite set of closed-form solutions $\{\prm{\Vtheta_0^*}, \prm{\Vtheta_1^*},
\ldots, \prm{\Vtheta_r^*}\}$. The global maximum can then be identified by simply
evaluating and comparing the log-likelihood at each stationary point. This is the ideal
scenario, realised for example in Gaussian models and in certain exponential family
distributions.

When closed-form solutions are unavailable, one must resort to numerical optimisation. If
the likelihood $L(\vrt|\vlx)$ can be evaluated for any $\Vtheta$, its gradient can also be
computed---either analytically or through automatic differentiation---and \alert{gradient
ascent} methods can be applied iteratively:
\begin{myequation}
	\prm{\Vtheta^{(i+1)}}=\prm{\Vtheta^{(i)}}+\eta\gradient{\vrt}{L(\vrt|\vlx)}\big\vert_{\prm{\Vtheta^{(i)}}},
\end{myequation}
where $\eta > 0$ is a step size (learning rate). Starting from some initial parameter value
$\prm{\Vtheta^{(0)}}$, this procedure follows the gradient of the likelihood uphill and
converges to a stationary point, which in general will be a local maximum. The quality of
the solution therefore depends on both the initialisation and the geometry of the likelihood
surface.

\subsection*{Maximum Likelihood from a Dataset}

In practice, parameter values are determined not from a single observation but from a
whole dataset $\X = \{\vlxu, \ldots, \vlxn\}$ of $n$ observations. Each observation is
assumed to be an independent realisation of the same random variable $\x$ distributed
according to $p(\x;\Vtheta)$---the standard \emph{i.i.d.} (independent and identically
distributed) assumption.

The independence assumption has a crucial consequence: the joint probability of the entire
dataset factorises as a product of individual observation probabilities. The likelihood of the
dataset is therefore:
\begin{myequation}
L(\vrt|\X)\defeq p(\X;\vrt)=\prod_{i=1}^n p(\vlxi ;\vrt).
\end{myequation}

\begin{figure}[ht]
\begin{center}
\begin{tikzpicture}[latent/.style={circle,draw,fill=latent},
visible/.style={circle,outer sep=0pt,draw,fill=observed},
param/.style={circle,draw, fill=param},
plate/.style={rounded corners,minimum height=1.5cm,minimum width=1.5cm,draw,fill=box},
lbl/.style={circle}]
\draw[line width=.5pt]
(2,0) node[param] (B) {$\Vtheta$}
(0,0) node[rounded corners,minimum height=2.1cm,minimum width=1.9cm,draw,fill=box] {}
(0,0) node[visible] (C) {$\vlxi$}
(.7,-.8) node[lbl] {$n$}
;
\draw[->,>=stealth',shorten >=1pt] (B) -- (C);
\end{tikzpicture}
\caption{The dataset consists of a set of $n$ independent realisations of the random
variable $\x$ distributed according to a given parametric distribution $p(\x;\Vtheta)$ with
parameter $\Vtheta$.}\label{bn1}
\end{center}
\end{figure}

The graphical model in Figure~\ref{bn1} represents this structure compactly: the plate
notation indicates that the $n$ observations $\vlxi$ are conditionally independent given
$\Vtheta$, all sharing the same generative mechanism parameterised by $\Vtheta$.

The best parameter value $\prts$ to model the dataset $\vlX$ is found by maximising either
the likelihood or, more conveniently, the log-likelihood of the entire dataset:
\begin{myequation}
	L(\vrt|\vlX)=\prod_{i=1}^n p(\vlxi;\Vtheta),
\end{myequation}
\begin{myequation}
	l(\vrt|\vlX)=\sum_{i=1}^n \log p(\vlxi;\Vtheta).
\end{myequation}
The log-likelihood of a dataset is thus a sum of individual log-likelihoods, one per data
point. This additive structure is computationally convenient and conceptually natural: each
observation contributes independently to the total evidence about the parameters. All the
considerations developed above for the single-observation case apply directly, with
$L(\vrt|\vlx)$ replaced by $L(\vrt|\vlX)$. The single-observation setting is simply the
special instance $n=1$.

\section*{Latent variable models}\label{sec:latentmodels}
\subsection*{Beyond Point Estimation: Latent Variable Models}

Inferring $\prts$ by maximum likelihood provides a useful but limited description of the
data.\footnote{In the following we will refer to a set $\vlX$ of observations, which in a
particular case shall be a single observation $\vlx$.} Maximum likelihood estimation tells
us which member of the parametric family $\{p(\x;\vrt)\}$ best fits the observations, but it
does not reveal the internal structure or organisation of the data. In unsupervised learning,
however, we are often interested in discovering patterns, clusters, or hidden dependencies
within a set of observations $\vlX$---questions that go beyond what a simple parametric
density estimate can answer.

To address these richer goals, more sophisticated probabilistic models can be introduced.
These models augment the observed variables $\x$ with additional \emph{latent} random
variables $\z = (z_1, \ldots, z_q)$ defined on a domain $\mZ$, which are never directly
observed but whose inferred values carry the structural information we wish to extract from
the data. The observed distribution $p(\x)$ is then expressed in terms of this richer joint
model by marginalising over the latent variables:
\begin{myequation}
p(\x) = \int_{\mZ} p(\x, \z)\, d\z = \int_{\mZ} p(\x|\z)\, p(\z)\, d\z.
\end{myequation}
This factorisation decomposes the model into two interpretable components: the
\emph{prior} $p(\z)$ over the latent variables, which encodes our beliefs about the hidden
structure before observing any data, and the \emph{conditional} $p(\x|\z)$, which describes
how the observed data is generated from a given latent configuration. By choosing
expressive forms for these two components, latent variable models can represent complex,
multimodal distributions over $\x$ that would be difficult to capture with a simple
parametric family on $\x$ alone. The challenge, as we will see, is that this increased
expressiveness comes at the cost of a more difficult inference and learning problem.

More precisely, we consider a probabilistic model $p(\x,\z; \Vtheta)$ where $\Vtheta=
\Vphi\cup\Vpsi$, involving an \alert{observed} random variable $\x$---whose instantiations
can be directly observed---and a \alert{latent} random variable $\z$, whose instantiations
are never accessible to the learner. The variable $\z$ is called \emph{latent} (from the
Latin for ``hidden'') precisely because it exists in the model only as a theoretical construct
that explains the variability in $\x$. Intuitively, one can think of $\z$ as encoding an
underlying cause, a cluster membership, or a compact representation of the observed data
point $\x$.

In the general case, both the domain $\mathcal{X}$ of $\x$ and the domain $\mathcal{Z}$
of $\z$ are continuous. In the simpler discrete case, where $\mathcal{Z}$ is a finite or
countable set, all integrals over $\mathcal{Z}$ are replaced by sums; the conceptual
structure of the model remains entirely the same, and the clustering models introduced in
the previous chapter---such as Gaussian mixture models---belong precisely to this discrete-
latent-variable regime.

The joint distribution factorises as $p(\x,\z; \Vtheta)=p(\x|\z;\Vtheta)p(\z;\Vtheta)$, and
we assume that the parameter vector $\Vtheta$ can be partitioned into two disjoint subsets
$\Vphi$ and $\Vpsi$, where $\Vphi$ governs the \emph{prior distribution} $p(\z;\Vphi)$
of the latent variable, while $\Vpsi$ governs the \emph{conditional distribution}
$p(\x|\z;\Vpsi)$ of the observed variable given the latent one. This separation reflects the
generative structure of the model: first, a latent value is drawn from the prior; then, an
observation is generated according to the conditional distribution at that latent value. We
may therefore write:
\begin{myequation}
p(\x,\z; \Vtheta)=p(\x|\z;\Vpsi)p(\z;\Vphi).
\end{myequation}

This factorisation is clean and intuitive, but it conceals a fundamental computational
difficulty. Even if we assume that the joint distribution $p(\x,\z;\Vtheta)$ is
tractable---meaning its value and its gradient with respect to $\Vtheta$ can be computed
efficiently---the \emph{marginal} distribution of the observed variable, known as the
\emph{evidence}, requires integrating out the latent variable:
\begin{myequation}
p(\x;\Vtheta)=\int_{\mZ}p(\x,\dz; \Vtheta)d\dz=\int_{\mZ}p(\x|\dz;\Vpsi)p(\z;\Vphi)d\dz.
\end{myequation}
This integral is, in general, intractable. The domain $\mathcal{Z}$ is typically
high-dimensional, rendering both analytical solutions and numerical approximations
computationally prohibitive: any quadrature-based approach would require evaluating the
integrand at an exponentially large number of points. This intractability of the evidence is
the central computational challenge in latent variable models, and it motivates the
approximate inference methods we will develop in the remainder of this chapter.

The same difficulty propagates to the posterior distribution of the latent variable given the
observation. By Bayes' theorem:
\begin{myequation}
p(\z|\x;\Vtheta)=\frac{p(\x,\z;\Vtheta)}{p(\x;\Vtheta)}.
\end{myequation}
Computing this posterior requires evaluating the evidence $p(\x;\Vtheta)$ in the
denominator, which is exactly what we cannot compute. The posterior $p(\z|\x;\Vtheta)$ is
of great practical importance---it tells us, given an observed data point, what the
distribution over latent representations looks like---but it is generally inaccessible in
closed form. To make progress, we need to understand precisely where the boundary
between tractable and intractable lies, and it is to this end that we now introduce two
working hypotheses.

To make the discussion precise, we introduce two hypotheses that delineate the boundary
between tractable and intractable regimes in latent variable models. These are not merely
technical conveniences---they reflect a fundamental asymmetry at the heart of the
challenge: the joint distribution over observed and latent variables is often manageable,
while the marginal distribution over observed variables alone is typically not.

The first hypothesis establishes a baseline level of tractability that we will assume
throughout the entire development.

\begin{description}
\item[Hypothesis 1]
For all values $\vlx, \lz, \prt$, it is feasible to compute the joint probability
$p(\vlx,\lz; \prt)$. Moreover, it is also feasible to find local maxima with respect to $\vrt$
of such distribution, or of its logarithm, either because the stationarity equation
$\gradient{\vrt}{p(\vlx,\lz; \vrt)}=\Zero$ admits closed-form solutions, or because the
gradient itself can be evaluated at any $\prt$, making it possible to apply numerical
optimisation methods such as gradient ascent. Under this hypothesis, it is therefore
possible to compute a (local) maximum $\prts$ of the complete-data likelihood:
\begin{myequation}
\prts=\argmax{\vrt}{p(\vlx,\lz; \vrt)}=\argmax{\vrt}{\log p(\vlx,\lz; \vrt)}=\argmax{\vrt}{l(\vrt|\vlx,\lz)}.
\end{myequation}
\end{description}

Hypothesis 1 is a baseline assumption that we will maintain throughout. It states that the
\emph{complete-data} likelihood---the likelihood of the joint observation $(\vlx, \lz)$---is
computationally accessible. This is a reasonable assumption for most latent variable
models of practical interest, since the joint distribution $p(\vlx,\lz;\prt) =
p(\vlx|\lz;\Vpsi)\,p(\lz;\Vphi)$ typically factorises into two components that are
individually simple and tractable. The difficulty, as we have seen, arises not from the
joint but from the marginal.

This leads to the second hypothesis, which addresses precisely whether the marginal
likelihood over the observed variables is tractable. Unlike Hypothesis 1, this condition may
or may not hold depending on the specific model, and the two cases give rise to
qualitatively different algorithmic strategies.

\begin{description}
\item[Hypothesis 2]
It is feasible, for all $\vlx$ and $\prt$, to compute the marginal probability
\begin{myequation}
p(\vlx; \prt)=\int_{\mZ}p(\vlx, \dz;\prt)\,d\dz.
\end{myequation}
As a consequence, the gradient $\gradient{\vrt}{p(\vlx; \vrt)}|_{\vrt=\prt}$ can be
evaluated, either in closed form or by automatic differentiation, and finding local maxima
of the observed-data likelihood is feasible, at least via gradient ascent. Equivalently,
when this hypothesis holds it is also feasible to compute the posterior distribution of the
latent variable given the observation:
\begin{myequation}
p(\z|\vlx;\prt)=\frac{p(\vlx,\z;\prt)}{p(\vlx;\prt)},
\end{myequation}
since the denominator is now available. This means that both the model parameters
$\theta_i$ and the latent values $\lat{\z_i}$ associated with each observation can in
principle be inferred: the parameters through maximum likelihood estimation, and the
latent variables by reasoning probabilistically under the posterior $p(\z|\vlx;\prt)$.
\end{description}

When Hypothesis 2 holds, computing the marginal probability $p(\vlx;\prt)$ is feasible for
any observation $\vlx$ and any parameter value $\prt$. This single fact is sufficient to
reduce the entire learning and inference problem to a standard maximum likelihood
estimation problem. Recall that in standard MLE, the only requirement is that one can
evaluate the likelihood $L(\vrt|\vlx) = p(\vlx;\vrt)$---and optionally its gradient---at any
candidate parameter value $\vrt$. Hypothesis 2 guarantees exactly this, and there is
therefore no obstacle to applying the same optimisation machinery one would use in a
model without latent variables at all. In other words, the presence of latent variables in
the model is, under Hypothesis 2, merely a modelling choice that enriches the
expressiveness of the distribution over $\x$---it does not introduce any additional
algorithmic difficulty. The latent variables have been effectively marginalised out, and what
remains is a perfectly ordinary parametric density $p(\vlx;\prt)$ that can be evaluated
pointwise.

The posterior $p(\z|\vlx;\prt)$ is also available as a by-product, since it requires only
dividing the tractable joint $p(\vlx,\z;\prt)$ by the now-tractable marginal $p(\vlx;\prt)$:
\begin{myequation}
p(\z|\vlx;\prt) = \frac{p(\vlx,\z;\prt)}{p(\vlx;\prt)}.
\end{myequation}
Since both quantities in this ratio are computable, the posterior is computable as well. This
means that once the parameters have been estimated via MLE, one can also perform
inference on the latent variables---assigning to each observation $\vlx$ a full probability
distribution over the latent states that generated it.

This favourable situation is, however, the exception rather than the rule. In the vast
majority of practically useful latent variable models---including mixture models, factor
analysis, and deep generative models---the marginal likelihood
\begin{myequation}
p(\vlx;\prt) = \int_{\mZ} p(\vlx, \dz;\prt)\, d\dz
\end{myequation}
involves a high-dimensional integral that is analytically intractable and numerically
infeasible to approximate. It is precisely in this regime that specialised algorithms become
necessary. The \alert{Expectation-Maximisation} (EM) algorithm, which we present shortly,
is the principal and most widely used tool for parameter estimation in this setting. Its central
insight is to circumvent the intractability of the marginal likelihood by constructing and
optimising a tractable surrogate---a lower bound on the log-likelihood built from the
complete-data likelihood which, by Hypothesis 1, remains accessible throughout.

\subsection*{Generation in Latent Variable Models}

Once a latent variable model $p(\x,\z;\Vtheta)$ has been specified and its parameters
estimated, it can be used not only for inference but also for \emph{generation}---that is,
for producing new synthetic observations in $\mathcal{X}$ that resemble those drawn from
the true data distribution. The generative capability of latent variable models is one of their
most compelling properties, and understanding it requires distinguishing between two
qualitatively different generation policies.

The most natural generation procedure follows directly from the probabilistic structure of
the model. Since the joint distribution factorises as $p(\x,\z;\Vtheta) =
p(\x|\z;\Vpsi)\,p(\z;\Vphi)$, new observations can be produced by ancestral sampling: one
first draws a latent value $\lz$ from the prior distribution $p(\z)$, and then draws an
observation $\val\x$ from the conditional distribution $p(\x|\lz)$, as illustrated in
Figure~\ref{f:gen}.\footnote{We shall not consider here the problem of how to sample
efficiently from arbitrary distributions.} This two-stage procedure has a clean and intuitive
interpretation. The latent variable $\lz$ acts as a selector: different values of $\lz$ index
different ``modes'' or ``components'' of the data distribution, each corresponding to a
distinct region or cluster in the observation space $\mathcal{X}$. Given $\lz$, the
conditional distribution $p(\x|\lz)$ then generates a specific observation consistent with
that mode. It is precisely this hierarchical structure---a continuous or discrete space of
latent configurations, each associated with a local distribution over observations---that
gives latent variable models their flexibility to represent complex, multimodal distributions.

\begin{figure}[ht]
\begin{center}
  \begin{tikzpicture}[>=latex']
	  \fill [elat!10] (8,7) rectangle (13,11);
    \node[elat] (c1) at (8.5,10.5) {$\mathcal Z$};

	  \draw [fill=elat!30, thick, draw=elat, densely dashed]
      plot [smooth cycle, tension=1] coordinates {(10,8.6) (12.4,8) (11,10)};
    \draw [fill=elat!30, draw=elat, densely dashed]
      plot [smooth cycle, tension=1] coordinates {(10.2,8.8) (12.2,8.2) (11.1,9.7)};
    \draw [fill=elat!30, draw=elat, densely dashed]
      plot [smooth cycle, tension=1] coordinates {(10.4,9) (12,8.5) (11.1,9.6)};
    \draw [fill=elat!30, draw=elat, densely dashed]
      plot [smooth cycle, tension=1] coordinates {(10.5,9.2) (11.5,8.8) (11.1,9.4)};
    \node[elat] (z) at (11,10.3) {$p(\z)$};

    \fill [odata!10] (0,3) rectangle (7,11);
    \node[odata] (c3) at (0.5,10.5) {$\mathcal X$};

    \draw [fill=odata!30, thick,draw=odata, densely dashed]
      plot [smooth cycle, tension=1] coordinates {(.5,4) (2,3.6) (1.2,6) (.5,5.6)};
    \draw [fill=odata!30, draw=odata, densely dashed]
      plot [smooth cycle, tension=1] coordinates {(.7,4.1) (1.7,3.8) (1.2,5.6) (.8,5.4)};
    \draw [fill=odata!30, draw=odata, densely dashed]
      plot [smooth cycle, tension=1] coordinates {(1.1,3.9) (1.6,4) (1,5.2)};
    \node[odata] (x1) at (2.7,3.5) {$p(\x|\overline{\z}_1)$};

	  \draw [fill=odata!30, thick,draw=odata, densely dashed]
      plot [smooth cycle, tension=1] coordinates {(1.7,8.6) (3.9,8.8) (4.3,10) (1.7,9.6)};
    \draw [fill=odata!30, draw=odata, densely dashed]
      plot [smooth cycle, tension=1] coordinates {(3,8.6) (4,9.8) (2.2,9.6) (1.8,9)};
    \draw [fill=odata!30, draw=odata, densely dashed]
      plot [smooth cycle, tension=1] coordinates {(3,8.7) (3.8,9.6) (2.4,9.4) (2,9)};
    \node[odata] (x2) at (1.5,10) {$p(\x|\overline{\z}_2)$};

    \draw [fill=odata!30, thick,draw=odata, densely dashed]
      plot [smooth cycle, tension=1] coordinates {(3.3,6.8) (6.2,6.4) (4.6,8) (3.1,7.6)};
    \draw [fill=odata!30, draw=odata, densely dashed]
      plot [smooth cycle, tension=1] coordinates {(3.8,6.9) (6,6.6) (4.6,7.6) (3.4,7.5)};
    \draw [fill=odata!30, draw=odata, densely dashed]
      plot [smooth cycle, tension=1] coordinates {(4,7) (5.9,6.7) (4.6,7.4) (3.8,7.5)};
    \node[odata] (x3) at (3,6.4) {$p(\x|\overline{\z}_3)$};

    \coordinate (a) at (11.5, 8.5);
    \coordinate (b) at (1, 5);
    \coordinate (c) at (4, 9.5);
    \coordinate (d) at (11, 9.5);
    \coordinate (e) at (5, 7);
    \coordinate (f) at (10.5, 8.7);

   	\draw [->, line width=.5mm, color=xkcdDustyOrange] (a) to [bend left=20] (b);
	  \draw [->, line width=.5mm, color=xkcdDustyOrange] (d) to [bend right=20] (c);
	  \draw [->, line width=.5mm, color=xkcdDustyOrange] (f) to [bend right=10] (e);

    \fill[elat] (a) circle [radius=2pt] node [anchor=north] {$\overline{\z}_1$};
    \fill[elat] (d) circle [radius=2pt] node [anchor=north] {$\overline{\z}_2$};
    \fill[elat] (f) circle [radius=2pt] node [anchor=north] {$\overline{\z}_3$};

    \fill[odata] (b) circle [radius=2pt] node [anchor=north] {$\vlxu$};
    \fill[odata] (c) circle [radius=2pt] node [anchor=north] {$\ox_2$};
    \fill[odata] (e) circle [radius=2pt] node [anchor=north] {$\ox_3$};
  \end{tikzpicture}
\caption{Generation of new items in a latent variable model: each latent value $\overline{\z}_i$
sampled from the prior $p(\z)$ determines a local conditional distribution
$p(\x|\overline{\z}_i)$ over the observation space, from which a new data point is then
sampled.}\label{f:gen}
\end{center}
\end{figure}

A second, qualitatively different generation policy arises when one wishes to produce a
new observation $\val{\overline\x} \in \mathcal{X}$ that is related to---or in many practical
contexts, similar to---a given reference observation $\val\x \in \mathcal{X}$. Rather than
sampling the latent variable freely from the prior, one first \emph{encodes} the reference
observation by sampling from the posterior distribution $p(\z|\val\x)$, obtaining a latent
value $\lz$ that is informed by $\val\x$ and captures its essential structure. One then
\emph{decodes} this latent value by sampling the conditional distribution $p(\x|\lz)$,
producing the new observation $\val{\overline\x}$. This encode-then-decode procedure
introduces a form of structured generation: the new sample is not drawn blindly from the
model, but is anchored to the reference observation through the shared latent
representation. The resulting $\val{\overline\x}$ can be thought of as a plausible variation
or reconstruction of $\val\x$, consistent with the patterns the model has learned from data.
This mechanism is precisely the generative procedure underlying Variational Autoencoders
and related models, where an encoder network approximates the posterior $p(\z|\x)$ and a
decoder network parameterises the conditional $p(\x|\z)$, enabling both efficient inference
and flexible generation.

\subsection*{Inference in Latent Variable Models}

Having described the generative structure of latent variable models, we now turn to the
inference problem: given observed data, how do we estimate the model parameters and
reason about the latent variables? This is where the computational difficulties described
above become central, and where the gap between the complete-data and observed-data
likelihoods takes on its full practical significance.

The inference setting involves two random variables: an \alert{observed random variable}
$\x$, whose instantiation $\ox$ is available to the learner, and a \alert{latent variable}
$\z$, whose instantiation $\lz$ is never directly observable. Together, $(\ox, \lz)$
constitute the \alert{complete dataset}---the full information that would be needed to
evaluate the complete-data likelihood. In practice, of course, only $\ox$ is available, and
$\lz$ must be reasoned about indirectly through the model.

In the general case, we work with a dataset $\X = \{\vlxu, \ldots, \vlxn\}$ of $n$
observations, each regarded as an independent instantiation of $\x$ generated according to
the latent variable model. This means we postulate the existence of a corresponding set
$\Z = \{\lat{\z_1}, \ldots, \lat{\z_n}\}$ of latent instantiations, one for each observation,
where each $\lat{\z_i}$ was sampled from the prior $p(\z;\Vphi)$ and each $\vlxi$ was
then sampled from the conditional $p(\x|\lat{\z_i};\Vpsi)$. The latent values $\lat{\z_i}$
are never observed---their existence is a modelling assumption that explains how the data
was generated, but they remain hidden throughout.

The independence structure of this generative process has immediate and important
consequences for the dataset-level distributions. Since each pair $(\vlxi, \lat{\z_i})$ is
generated independently of all others, all joint, marginal, and conditional distributions over
the full dataset factorise into products of the corresponding single-observation
distributions:
\begin{myalign}
p(\Z)&=\prod_{i=1}^n p(\lat{\z_i}),\\
p(\X|\Z)&=\prod_{i=1}^n p(\vlxi|\lat{\z_i}),\\
p(\X,\Z)&=\prod_{i=1}^n p(\vlxi,\lat{\z_i}),\\
p(\X)&=\prod_{i=1}^n p(\vlxi),\\
p(\Z|\X)&=\prod_{i=1}^n p(\lat{\z_i}|\vlxi).
\end{myalign}
These product forms are a direct consequence of the conditional independence assumptions
encoded in the model: given their respective latent variables, the observations are mutually
independent. This structure is represented compactly in the graphical model of
Figure~\ref{bn2}, where the plate notation indicates that the $n$ pairs $(\lat{\z_i},
\vlxi)$ are independent replications of the same generative process, governed by shared
parameters $\prm\Vphi$ and $\prm\Vpsi$.

\begin{figure}[ht]
\begin{center}
\begin{tikzpicture}[
latent/.style={circle,draw,fill=latent!40!white},
visible/.style={circle,outer sep=0pt,draw,fill=observed},
param/.style={circle,draw, fill=param},
plate/.style={rounded corners,minimum height=1.5cm,minimum width=1.5cm,draw,fill=box},
lbl/.style={circle}]
\draw[line width=.5pt]
(2,1.5) node[param] (D) {$\prm\Vphi$}
(2,0) node[param] (B) {$\prm\Vpsi$}
(0,.7) node[rounded corners,minimum height=3.1cm,minimum width=1.9cm,draw,fill=node!50] {}
(0,1.5) node[latent] (A) {$\lat{\z_i}$}
(0,0) node[visible] (C) {$\vlxi$}
(.7,-.6) node[lbl] {\footnotesize $n$};
\draw[->,>=stealth',shorten >=1pt] (B) -- (C);
\draw[->,>=stealth',shorten >=1pt] (A) -- (C);
\draw[->,>=stealth',shorten >=1pt] (D) -- (A);
\end{tikzpicture}
\caption{The dataset is composed of a set of $n$ independent realisations of a pair of
random variables $\x, \z$, with the latent variable $\z$ distributed according to
$p(\z;\prm\Vphi)$ and the observed variable $\x$ dependent on $\z$ and distributed
according to the conditional $p(\x|\z;\prm\Vpsi)$.}\label{bn2}
\end{center}
\end{figure}

In the following, we work with the general case of a dataset of $n$ pairs
$\{(\vlxu,\lat{\z_1}),\ldots,(\vlxn,\lat{\z_n})\}$, keeping in mind that the
single-observation case $(\ox,\lz)$ is simply the special instance $n=1$ and that all
results extend straightforwardly between the two settings.

Since each observation $\vlxi$ is available while the corresponding latent value
$\lat{\z_i}$ is not, our knowledge of the latent variable is necessarily probabilistic. The
appropriate object for reasoning about $\lat{\z_i}$ given $\vlxi$ is the posterior
distribution $p(\z|\vlx;\prt)$, which by Bayes' theorem takes the form:
\begin{myalign}
p(\lat{\z_i}|\vlxi;\prt)&=\frac{p(\vlxi|\lat{\z_i};\prm\Vpsi)\,p(\lat{\z_i};\prm\Vphi)}{p(\vlxi;\prm\Vtheta)}=\frac{p(\vlxi|\lat{\z_i};\prm\Vpsi)\,p(\lat{\z_i};\prm\Vphi)}{\int_{\mathcal{Z}}p(\vlxi,\z;\prt)\,d\z}.
\end{myalign}
The numerator is tractable by Hypothesis 1: it is simply the joint probability
$p(\vlxi,\lat{\z_i};\prt)$, which we assume can be evaluated efficiently. The denominator,
however, is the evidence $p(\vlxi;\prm\Vtheta)$, obtained by integrating the joint over all
possible latent values---precisely the intractable integral we have been grappling with
throughout. We therefore proceed under the assumption that Hypothesis 2 does not hold:
the evidence $p(\vlxi;\prm\Vtheta)$ cannot be feasibly computed, which in turn makes the
posterior $p(\lat{\z_i}|\vlxi;\prm\Vtheta)$ inaccessible by direct calculation. Both
parameter estimation and latent variable inference are thus blocked by the same
fundamental obstacle, and it is to overcome this obstacle that we now develop the EM
algorithm.

\paragraph{The gap between complete and observed likelihoods.}

The inference problem identified above has a direct and critical consequence for parameter
estimation. Following the maximum likelihood principle, our objective is to maximise the
\emph{observed-data} log-likelihood:
\begin{myalign}
l(\vrt|\vlX)&=\log p(\X;\Vtheta)=\sum_{i=1}^{n}\log p(\vlxi;\Vtheta),
\end{myalign}
since $\vlX$ is the only dataset actually available to us. Each term in this sum requires
evaluating $\log p(\vlxi;\Vtheta)$---the log-marginal probability of $\vlxi$ under the
model---which in turn requires integrating out the latent variable. Under the assumption
that Hypothesis 2 fails, this is intractable: we cannot even evaluate the objective function
we wish to optimise, let alone compute its gradient.

The situation would be entirely different if the latent variables were observed. If the
complete dataset $(\vlX, \Z)$ were available, we could work with the
\emph{complete-data} log-likelihood:
\begin{myalign}
l(\vrt|\vlX, \Z)&=\log p(\X,\Z;\Vtheta)=\sum_{i=1}^{n}\log p(\vlxi , \lat{\z_i};\vrt).
\end{myalign}
Each term in this sum involves only the joint probability $p(\vlxi, \lat{\z_i};\vrt)$, which
by Hypothesis 1 is tractable to compute and optimise. The complete-data log-likelihood is
therefore a well-behaved objective---it is accessible, differentiable, and amenable to
standard maximum likelihood techniques. The fundamental difficulty is simply that $\Z$ is
not available, so this objective cannot be evaluated either.

We are thus caught between two likelihoods: the observed-data likelihood, which
corresponds to the right objective but is intractable, and the complete-data likelihood,
which is tractable but depends on unobserved quantities. As we shall see, the
\alert{Expectation-Maximisation} (EM) algorithm resolves this tension elegantly by
constructing a surrogate objective that bridges the two. Rather than maximising the
complete-data log-likelihood directly---which would require knowledge of $\Z$---the EM
algorithm replaces the unknown latent variables with their expected values under the
posterior distribution $p(\Z|\vlX;\Vtheta)$ computed at the current parameter estimate.
This yields a tractable lower bound on the observed-data log-likelihood that can be
maximised with respect to the parameters, and the procedure is iterated to convergence.
The precise construction and justification of this approach is the subject of the next section.

\section*{The Evidence Lower Bound}

Having established that direct maximisation of the observed-data log-likelihood
$\log p(\ox;\Vtheta)$ is generally intractable, we now develop the key mathematical tool
that makes approximate maximisation possible: the \alert{Evidence Lower Bound} (ELBO).
The central idea is to introduce an auxiliary distribution $q(\z)$ over the latent variable and
use it to construct a tractable lower bound on the log-likelihood. By iteratively tightening
and maximising this lower bound, we will obtain a principled algorithm for parameter
estimation in latent variable models.

\subsection*{Derivation of the ELBO}

Let us first consider the case of a single observation $\ox$, and let $q(\z)$ be any
probability distribution over the latent variable. The choice of $q(\z)$ is, for the moment,
completely arbitrary---it can be any valid probability distribution on $\mathcal{Z}$, and the
derivation that follows holds regardless of this choice. The key insight is that we can
introduce $q(\z)$ into the expression for $\log p(\ox;\prt)$ without changing its value, by
exploiting the fact that any probability distribution integrates to one. This manipulation,
while seemingly circular, will allow us to decompose the log-likelihood into two
interpretable and useful terms, each with a clear probabilistic meaning.

For any $\prt$, the following chain of equalities holds:
\begin{myalign}
\log p(\ox ;\prt)&=\log p(\ox ;\prt){\int_{\mZ}q(\dz)d\dz} \\
&=\int_{\mZ}q(\dz)\log p(\ox ;\prt)d\dz \\
&=\int_{\mZ} q(\dz)\log\frac{p(\ox,\dz;\prt)}{p(\dz|\ox;\prt)}d\dz \\
&=\int_{\mZ} q(\dz)\left(\log p(\ox,\dz;\prt)-\log  p(\dz|\ox;\prt)\right)d\dz \\
&=\int_{\mZ} q(\dz)\left(\log\frac{p(\ox,\dz;\prt)}{ q(\dz)}-\log\frac{p(\dz|\ox;\prt)}{ q(\dz)}\right)d\dz \\
&=\int_{\mZ}q(\dz)\log\frac{p(\ox,\dz;\prt)}{q(\dz)}d\dz-\int_{\mZ}q(\dz)\log\frac{p(\dz|\ox;\prt)}{q(\dz)}d\dz \\
&=\mL(q,\ox, \prt)+\mK(q,\ox, \prt).
\end{myalign}

It is worth tracing through the logic of each step carefully, as the derivation is compact but
non-obvious. The first two steps simply use the fact that $q(\z)$ integrates to one, moving
$\log p(\ox;\prt)$---which does not depend on $\z$---inside the integral. The third step
rewrites $p(\ox;\prt)$ using the definition of conditional probability: since
$p(\ox,\z;\prt) = p(\z|\ox;\prt)\,p(\ox;\prt)$, we have $p(\ox;\prt) =
p(\ox,\z;\prt)/p(\z|\ox;\prt)$. The fifth step introduces $q(\z)$ symmetrically in both
terms by adding and subtracting $\log q(\z)$---a manipulation that leaves the value
unchanged since $q(\z)/q(\z) = 1$. Finally, the integral separates into two terms, each
taking the form of an expectation under $q(\z)$ and admitting a clear probabilistic
interpretation.

The two resulting quantities are defined as:
\begin{myalign}
\mL(q,\ox, \prt)& \defeq\ev{\dz\sim q(\dz)}{\log\frac{p(\ox,\dz;\prt)}{q(\dz)}}=-\ev{\dz\sim q(\dz)}{\log\frac{q(\dz)}{p(\ox,\dz;\prt)}}=-\dkl{q(\z)}{p(\ox ,\z;\prt)},\\
\mK(q,\ox, \prt)&\defeq-\ev{\dz\sim q(\dz)}{\log\frac{p(\dz|\ox;\prt)}{q(\dz)}}=\ev{\dz\sim q(\dz)}{\log\frac{q(\dz)}{p(\dz|\ox;\prt)}}=\dkl{q(\z)}{p(\z|\ox ;\prt)},
\end{myalign}
where the \alert{Kullback-Leibler divergence} between two distributions $p_1(\z)$ and
$p_2(\z)$ is defined as:
\begin{myequation}
\dkl{p_1(\z)}{p_2(\z)}\defeq \ev{p_1(\z)}{\log\frac{p_1(\z)}{p_2(\z)}}=\int_{\mZ}p_1(\z)\log\frac{p_1(\z)}{p_2(\z)}d\z.
\end{myequation}
The KL divergence is a fundamental quantity in information theory and statistics that
measures the dissimilarity between two probability distributions: it quantifies how much
information is lost when $p_2$ is used to approximate $p_1$, or equivalently, how
surprised one would be on average if observations drawn from $p_1$ were evaluated under
$p_2$. When $\z$ is discrete, the integral is replaced by a sum over $\mathcal{Z}$.
Importantly, the KL divergence is \emph{not} symmetric in general---$\dkl{p_1}{p_2} \neq
\dkl{p_2}{p_1}$---and this asymmetry will have meaningful consequences for how we use it.

The two terms $\mL$ and $\mK$ have distinct and complementary interpretations. The
term $\mK(q,\ox, \prt)$ is a genuine KL divergence between $q(\z)$ and the true posterior
$p(\z|\ox;\prt)$: it measures precisely how far the auxiliary distribution $q$ is from the
posterior that we would ideally use to reason about the latent variable given the observation.
The term $\mL(q,\ox, \prt)$, on the other hand, is written in the form of a negative KL
divergence between $q(\z)$ and the joint distribution $p(\ox,\z;\prt)$. A technical
clarification is in order here: the standard KL divergence $\dkl{p_1(\z)}{p_2(\z)}$ is
defined between two probability distributions over the same space, meaning both $p_1$ and
$p_2$ should be normalised over $\z$. However, $p(\ox,\z;\prt)$ is a joint distribution over
the product space $\mathcal{X} \times \mathcal{Z}$, not a distribution over $\mathcal{Z}$
alone. Viewed as a function of $\z$ with $\ox$ fixed, it is generally not normalised:
integrating it over $\z$ gives $p(\ox;\prt)$, which is not $1$ in general. The expression
\begin{myequation}
\mL(q,\ox,\prt) = \int_{\mZ} q(\z) \log \frac{p(\ox,\z;\prt)}{q(\z)}\, d\z
\end{myequation}
therefore does not fit the strict definition of a negative KL divergence, because the
numerator in the logarithm is an unnormalised function of $\z$. The notation
$-\dkl{q(\z)}{p(\ox,\z;\prt)}$ is nonetheless widely adopted in the literature as a
deliberate and convenient abuse: the mathematical form is identical to that of a KL
divergence, and the same intuition applies---$\mL$ is large when $q(\z)$ places mass in
regions where $p(\ox,\z;\prt)$ is large, and small otherwise.

Before proceeding, it is worth pausing to understand the mathematical nature of both
$\mL$ and $\mK$. Neither is a function of $\z$, since $\z$ has been marginalised out by
the integration. Both depend on the observation $\ox$ (treated here as fixed) and on the
parameter $\Vtheta$. Most importantly, both are \emph{functionals} of the auxiliary
distribution $q$---their inputs include an entire probability distribution rather than a
finite-dimensional vector. Although both $\mL$ and $\mK$ individually depend on $q$,
their sum equals $\log p(\ox;\prt)$, which is entirely independent of $q$. This conservation
property is the key to understanding why the ELBO is useful: varying $q$ redistributes
value between $\mL$ and $\mK$ but leaves their sum---the log-likelihood---unchanged.
Increasing $\mL$ by choosing a better $q$ necessarily decreases $\mK$ by the same
amount, tightening the approximation to the true log-likelihood.

\subsection*{The ELBO as a Lower Bound}

The decomposition $\log p(\ox;\prt) = \mL(q,\ox,\prt) + \mK(q,\ox,\prt)$ becomes truly
useful once we invoke the fundamental properties of the KL divergence. For any two
distributions $p_1$ and $p_2$:
\begin{myalign}
	\dkl{p_1(\z)}{p_2(\z)}&\geq 0,\\
	\dkl{p_1(\z)}{p_2(\z)}&=0\quad\text{if and only if}\quad p_1(\z)=p_2(\z)\quad\text{for all }\z.
\end{myalign}
The non-negativity of the KL divergence is a classical result in information theory. It
follows from Jensen's inequality applied to the strictly convex function $-\log$: since
$-\log$ is convex, the expected value of $-\log(p_2/p_1)$ under $p_1$ is at least $-\log$
of the expected value of $p_2/p_1$, which equals zero by normalisation. Equality holds if
and only if the ratio $p_2(\z)/p_1(\z)$ is constant almost everywhere, which for normalised
distributions forces $p_1 = p_2$.

This non-negativity has an immediate and important consequence for our decomposition.
Since $\mK(q,\x,\prt) = \dkl{q(\z)}{p(\z|\x;\prt)} \geq 0$, it follows directly that:
\begin{myequation}
\mL(q,\x, \prt)=\log p(\x ;\prt)-\mK(q,\x, \prt)\leq\log p(\x ;\prt),
\end{myequation}
and therefore, for any observation $\ox$, any parameter value $\prt$, and any auxiliary
distribution $q(\z)$ whatsoever:
\begin{myequation}
\mL(\fq,\x, \prt)\leq\log p(\x;\prt)\hspace{1cm}\text{ whatever the distribution }\fq.
\end{myequation}
The quantity $\mL(q,\ox,\prt)$ is therefore a lower bound on the log-likelihood
$\log p(\ox;\prt)$, holding \emph{universally}---it does not depend on any particular choice
of $q$, and no assumption on the model or the observation is required. We call
$\mL(q,\ox,\prt)$ the \alert{evidence lower bound}, or \alert{ELBO}, for parameters $\prt$
and observation $\ox$. The name reflects both the role of $p(\ox;\prt)$ as the ``evidence''
in Bayesian terminology and the fact that $\mL$ provides a lower bound on its logarithm.

The gap between the ELBO and the true log-likelihood is not arbitrary: it is exactly
$\mK(q,\ox,\prt) = \dkl{q(\z)}{p(\z|\ox;\prt)}$, the KL divergence between the auxiliary
distribution $q$ and the true posterior over the latent variable. This has a clear and
important interpretation: the ELBO is tight---that is, it equals the log-likelihood---if and
only if $q(\z) = p(\z|\ox;\prt)$ for all $\z$, i.e., when the auxiliary distribution coincides
with the true posterior. Any deviation of $q$ from the posterior introduces a gap, and the
size of that gap is precisely measured by the KL divergence between the two distributions.

Beyond its role as a bound, the ELBO admits an alternative decomposition that reveals
why it is computationally tractable:
\begin{myequation}
 \mL(q,\x, \prt)=\ev{\z\sim q(\z)}{\log p(\x,\dz;\prt)}-\ev{\z\sim q(\z)}{\log q(\dz)}.
\end{myequation}
The first term is the expected complete-data log-likelihood under the distribution $q(\z)$:
it averages $\log p(\x,\z;\prt)$ over the latent variable distributed according to $q$. By
Hypothesis 1, the complete-data log-likelihood is tractable to evaluate for any fixed $\z$,
and computing its expectation under a chosen $q$ is therefore feasible whenever
expectations under $q$ can be computed analytically or estimated by sampling. The second
term is the \emph{entropy} of $q(\z)$, defined as $-\ev{q(\z)}{\log q(\z)}$, which depends
only on $q$ and not on the model parameters $\Vtheta$. The entropy term acts as a
regulariser: it is large when $q$ is diffuse and spread out, and small when $q$ is
concentrated, thus penalising overly peaked choices of the auxiliary distribution.

Taken together, the ELBO differs from the expected complete-data log-likelihood only by
the entropy of $q(\z)$, which is a constant with respect to $\Vtheta$. This means that, for
a fixed $q$, maximising the ELBO over $\Vtheta$ is equivalent to maximising the expected
complete-data log-likelihood---a quantity that is tractable by Hypothesis 1. This is the
fundamental reason why the ELBO serves as a useful surrogate objective: it replaces the
intractable log-likelihood with a computable approximation whose tightness is controlled by
how well $q$ approximates the true posterior.

As illustrated in Figure~\ref{f:elbo}, different choices of the auxiliary distribution $q$
yield different lower bounds on the same log-likelihood $\log p(\x;\prt)$. Some choices
produce tighter bounds that closely track the true log-likelihood, while others produce
looser bounds. The figure also shows, for each $q_i$, the corresponding expected
complete-data log-likelihood $\ev{q_i}{\log p(\x,\z;\prt)}$ as a dashed curve: this lies
uniformly above the ELBO by the (constant) entropy of $q_i$, confirming the
decomposition above. The art of the EM algorithm, as we shall see, lies in choosing $q$
optimally at each iteration so as to make the bound as tight as possible before maximising
it over $\Vtheta$.

\begin{figure}[h]
\begin{center}
\begin{tikzpicture}[scale=1]
\begin{axis}[
    clip=false,axis lines=middle,
    xmin=-2, xmax=16, xlabel = $\x$,
    ymin=-3, ymax=1,  ylabel={},
    xticklabels = {},
    yticklabels = {},
    extra x ticks={},
    extra x tick style={grid=major,
        tick label style={anchor=north west}},
    extra x tick labels={},
    legend style={
        legend pos=outer north east,
    },]
\addplot [xkcdCadetBlue, line width=.4mm, smooth, tension=.9] coordinates {(-2, -2.9) (0,-2.7) (2,-1.6) (4,-1.4) (6,-1.3) (8,-2) (10,-.7) (12,-1.5) (14,-2.3)};
\addlegendentry{\footnotesize $\log p(\x;\prt)$}
\addplot [xkcdRed, smooth, tension=.9] coordinates {(-2, -3) (0,-3) (2,-2) (4,-1.5) (6,-1.8) (8,-2.5) (10,-2) (12,-2.2) (14,-2.7)};
\addlegendentry{\footnotesize $\mL(\func{q_1},\x, \prt)$}
\addplot [xkcdGreen,smooth, tension=.9] coordinates {(-2,-3.2) (0,-2.9) (2,-2) (4,-2.3) (6,-2.6) (8,-3) (10,-1) (12,-2) (14,-2.7)};
\addlegendentry{\footnotesize $\mL(\func{q_2},\x, \prt)$}
\addplot [xkcdOrange,smooth, tension=.9] coordinates {(-2,-3.2) (0,-3) (2,-2) (4,-2.1) (6,-2.3) (8,-3.2) (10,-1.8) (12,-2.1) (14,-3)};
\addlegendentry{\footnotesize $\mL(\func{q_3},\x, \prt)$}
\addplot [xkcdRed, smooth, dashed, tension=.9, yshift=2mm] coordinates {(-2, -3) (0,-3) (2,-2) (4,-1.5) (6,-1.8) (8,-2.5) (10,-2) (12,-2.2) (14,-2.7)};
\addlegendentry{\footnotesize $\ev{\func{q_1}}{\log p(\x,\dz;\prt)}$}
\addplot [xkcdGreen, smooth, dashed, tension=.9, yshift=1mm] coordinates {(-2,-3.2) (0,-2.9) (2,-2) (4,-2.3) (6,-2.6) (8,-3) (10,-1) (12,-2) (14,-2.7)};
\addlegendentry{\footnotesize $\ev{\func{q_2}}{\log p(\x,\dz;\prt)}$}
\addplot [xkcdOrange, smooth, dashed, tension=.9, yshift=1.7mm] coordinates {(-2,-3.2) (0,-3) (2,-2) (4,-2.1) (6,-2.3) (8,-3.2) (10,-1.8) (12,-2.1) (14,-3)};
\addlegendentry{\footnotesize $\ev{\func{q_3}}{\log p(\x,\dz;\prt)}$}
\end{axis}
\end{tikzpicture}
\end{center}
\caption{Each distribution $\func{q_i}$ provides a different lower bound (ELBO) on the
log-likelihood $\log p(\x;\prt)$, shown as the thick blue curve. The dashed curves show
the corresponding expected complete-data log-likelihoods $\ev{\func{q_i}}{\log
p(\x,\z;\prt)}$, which lie uniformly above the respective ELBOs by the entropy of
$\func{q_i}$.}\label{f:elbo}
\end{figure}

\subsection*{Iterative Maximisation via the ELBO}

The ELBO provides more than just a lower bound on the log-likelihood---it provides a
concrete strategy for maximising that log-likelihood iteratively, even when direct
maximisation is infeasible. The key observation is that the ELBO depends on two distinct
objects: the auxiliary distribution $q$, which controls how tight the bound is, and the
model parameters $\Vtheta$, which we ultimately wish to optimise. This suggests a natural
two-step iterative procedure that alternates between optimising over $q$ with $\Vtheta$
fixed, and optimising over $\Vtheta$ with $q$ fixed. Each step is individually tractable,
and together they drive the log-likelihood upward---a claim we will make precise below.

\paragraph{First step: tightening the bound.}
At iteration $k$, suppose we have a current parameter estimate $\prm{\Vtheta^{(k)}}$.
The log-likelihood $\log p(\ox;\prm{\Vtheta^{(k)}})$ is fixed at this point, but the ELBO
\begin{myequation}
\mL(q,\ox,\prm{\Vtheta^{(k)}})=\ev{q(\z)}{\log p(\ox,\dz;\prm{\Vtheta^{(k)}})}-\ev{q(\z)}{\log q(\dz)}\leq \log p(\ox;\prm{\Vtheta^{(k)}})
\end{myequation}
still depends on $q$, and different choices of $q$ yield different lower bounds on the same
log-likelihood value, as illustrated in Figure~\ref{f:elbo1}. We wish to find the distribution
$\func{q^{(k)}}$ that makes the bound as tight as possible, i.e., that maximises
$\mL(q,\ox,\prm{\Vtheta^{(k)}})$ over all valid distributions $q$. Since the log-likelihood
is fixed and the gap between the ELBO and the log-likelihood is exactly
$\mK(q,\ox,\prm{\Vtheta^{(k)}}) = \dkl{q(\z)}{p(\z|\ox;\prm{\Vtheta^{(k)}})}$,
maximising the ELBO over $q$ is equivalent to minimising this KL divergence. The
minimum is zero, achieved when $q(\z) = p(\z|\ox;\prm{\Vtheta^{(k)}})$, at which point the
ELBO touches the log-likelihood from below and the bound is exact. The optimal
distribution is therefore the true posterior of the latent variable under the current
parameter estimate:
\begin{myequation}
  \func{q^{(k)}}=\argmax{\vr{q}}{\mL(\vr{q},\ox, \prm{\Vtheta^{(k)}})} = p(\z|\ox;\prm{\Vtheta^{(k)}}).
\end{myequation}

\begin{figure}[h]
\begin{center}
\begin{tikzpicture}[scale=1]
\begin{axis}[
    clip=false,axis lines=middle,
    xmin=-2, xmax=16, xlabel = $\x$,
    ymin=-3, ymax=1,  ylabel={},
    xticklabels = {},
    yticklabels = {},
    extra x ticks={4},
    extra x tick style={grid=major,
        tick label style={anchor=north west}},
    extra x tick labels={$\ox$},
    legend style={
        legend pos=outer north east,
    },]
\addplot [xkcdCadetBlue, line width=.4mm, smooth, tension=.9] coordinates {(-2, -2.9) (0,-2.7) (2,-1.6) (4,-1.4) (6,-1.3) (8,-2) (10,-.7) (12,-1.5) (14,-2.3)};
\addlegendentry{\footnotesize $\log p(\x;\prm{\Vtheta^{(k)}})$}
\addplot [xkcdRed, smooth, tension=.9] coordinates {(-2, -3) (0,-3) (2,-2) (4,-1.5) (6,-1.8) (8,-2.5) (10,-2) (12,-2.2) (14,-2.7)};
\addlegendentry{\footnotesize $\mL(\func{q_1},\x, \prm{\Vtheta^{(k)}})$}
\addplot [xkcdGreen,smooth, tension=.9] coordinates {(-2,-3.2) (0,-2.9) (2,-2) (4,-2.3) (6,-2.6) (8,-3) (10,-1) (12,-2) (14,-2.7)};
\addlegendentry{\footnotesize $\mL(\func{q_2},\x, \prm{\Vtheta^{(k)}})$}
\addplot [xkcdOrange,smooth, tension=.9] coordinates {(-2,-3.2) (0,-3) (2,-2) (4,-2.1) (6,-2.3) (8,-3.2) (10,-1.8) (12,-2.1) (14,-3)};
\addlegendentry{\footnotesize $\mL(\func{q_3},\x, \prm{\Vtheta^{(k)}})$}
\filldraw (4,-1.5) [xkcdRed]circle (2pt)node{};
\filldraw (3.5,-1) node[above, font=\footnotesize] {$\mL(\func{q_1},\ox,\prm{\Vtheta^{(k)}})$};
\draw[-stealth]  (3.5,-1) -- (4,-1.5);
\end{axis}
\end{tikzpicture}
\end{center}
\caption{Different distributions provide different lower bounds to the log-likelihood at the
observation $\ox$, given the current parameter value $\prm{\Vtheta^{(k)}}$. The optimal
$q$ is the one whose ELBO is closest to the log-likelihood at $\ox$.}\label{f:elbo1}
\end{figure}

\paragraph{Second step: maximising the bound over parameters.}
Having fixed $\func{q^{(k)}}$, we now consider the ELBO as a function of $\Vtheta$ alone:
\begin{myequation}
\mL(\func{q^{(k)}},\ox,\Vtheta)=\ev{\func{q^{(k)}}(\z)}{\log p(\ox,\dz;\Vtheta)}-\ev{\func{q^{(k)}}(\z)}{\log\func{q^{(k)}}(\dz)}.
\end{myequation}
The second term is the entropy of $\func{q^{(k)}}(\z)$, written $\entr{}{\func{q^{(k)}}(\dz)}$,
which does not depend on $\Vtheta$ and is therefore a constant with respect to the
optimisation. Maximising the ELBO over $\Vtheta$ is thus equivalent to maximising only
the first term---the expected complete-data log-likelihood under $\func{q^{(k)}}$:
\begin{myequation}
Q(\func{q^{(k)}},\ox,\vrt) \defeq \ev{\func{q^{(k)}}(\z)}{\log p(\ox,\dz;\vrt)}.
\end{myequation}
This quantity is tractable by Hypothesis 1: the complete-data log-likelihood $\log
p(\ox,\z;\vrt)$ can be evaluated for any $\z$ and $\vrt$, and $\func{q^{(k)}}$ is a known
distribution, so computing its expectation is feasible. The new parameter estimate is then:
\begin{myequation}
  \prm{\Vtheta^{(k+1)}}=\argmax{\vrt}{\mL(\func{q^{(k)}},\ox,\vrt)}=\argmax{\vrt}{Q(\func{q^{(k)}},\ox,\vrt)}.
\end{myequation}

The geometric picture of this step is shown in Figure~\ref{f:elbo2}. With $\func{q^{(k)}}$
fixed, the ELBO $\mL(\func{q^{(k)}},\ox,\Vtheta)$ acts as a lower bound on the
log-likelihood $\log p(\ox;\Vtheta)$ for every value of $\Vtheta$---not just at the current
$\prm{\Vtheta^{(k)}}$, but everywhere. Maximising this lower bound over $\Vtheta$ yields
the new parameter estimate $\prm{\Vtheta^{(k+1)}}$, which achieves a higher value of
the ELBO than $\prm{\Vtheta^{(k)}}$. Since the ELBO is a lower bound on the
log-likelihood at every $\Vtheta$, the new log-likelihood $\log
p(\ox;\prm{\Vtheta^{(k+1)}})$ is at least as large as the new ELBO value---though, as we
discuss shortly, this does not immediately guarantee that the log-likelihood itself has
increased.

\begin{figure}[h]
\begin{center}
\begin{tikzpicture}[scale=1]
\begin{axis}[
    clip=false,axis lines=middle,
    xmin=-2, xmax=16, xlabel = $\x$,
    ymin=-3, ymax=1,  ylabel={},
    xticklabels = {},
    yticklabels = {},
    extra x ticks={4},
    extra x tick style={grid=major,
        tick label style={anchor=north west}},
    extra x tick labels={$\ox$},
    legend style={
        legend pos=outer north east,
    },]
\addplot [xkcdCadetBlue, line width=.4mm, smooth, tension=.9] coordinates {(-2, -2.9) (0,-2.7) (2,-1.6) (4,-1.4) (6,-1.3) (8,-2) (10,-.7) (12,-1.5) (14,-2.3)};
\addlegendentry{\footnotesize  $\log p(\x;\prm{\Vtheta^{(k)}})$}
\addplot [xkcdRed, smooth, tension=.9] coordinates {(-2, -3) (0,-3) (2,-2) (4,-1.5) (6,-1.8) (8,-2.5) (10,-2) (12,-2.2) (14,-2.7)};
\addlegendentry{\footnotesize  $\mL(\func{q^{(k)}},\x, \prm{\Vtheta^{(k)}})$}
\addplot [xkcdBrownyOrange,smooth, tension=.9] coordinates {(-2, -2.4) (0,-2.7) (2,-2.6) (4,-2.7) (6,-2.1) (8,-1.8) (10,-1.4) (12,-2.1) (14,-2.3)};
\addlegendentry{\footnotesize  $\mL(\func{q^{(k)}},\x, \prm{\Vtheta'})$}
\addplot [xkcdFlatGreen,smooth, tension=.9] coordinates {(-2,-2.7) (0,-2.4) (2,-1.7) (4,-.8) (6,-.6) (8,-1.5) (10,-2) (12,-1.7) (14,-2.5)};
\addlegendentry{\footnotesize  $\mL(\func{q^{(k)}},\x,\prm{\Vtheta^{(k+1)}})$}
\filldraw (4,-1.5) [xkcdRed]circle (2pt)node{};
\filldraw (-.5,-1.1) node[left, font=\footnotesize] {$\mL(\func{q^{(k)}},\ox,\prm{\Vtheta^{(k)}})$};
\draw[-stealth]  (-.5,-1.1) -- (4,-1.5);
\filldraw (4,-.8) [xkcdFlatGreen]circle (2pt)node{};
\filldraw (7.5,-.5) node[right, font=\footnotesize] {$\mL(\func{q^{(k)}},\ox,\prm{\Vtheta^{(k+1)}})$};
\draw[-stealth]  (7.5,-.5) -- (4,-.8);
\filldraw (4,-2.7) [xkcdBrownyOrange]circle (2pt)node{};
\filldraw (5,-3.2) node[right, font=\footnotesize] {$\mL(\func{q^{(k)}},\ox,\prm{\Vtheta'})$};
\draw[-stealth]  (5,-3.2) -- (4,-2.7);
\addplot [xkcdBarneyPurple, line width=.3mm, smooth, tension=.9] coordinates {(-2, -2.5) (0,-2.2) (2,-1.2) (4,-.7) (6,-.3) (8,-1.3) (10,-1.5) (12,-1.6) (14,-2.2)};
\addlegendentry{\footnotesize  $\log p(\x;\prm{\Vtheta^{(k+1)}})$}
\end{axis}
\end{tikzpicture}
\end{center}
\caption{With $\func{q^{(k)}}$ fixed, the ELBO $\mL(\func{q^{(k)}},\x,\Vtheta)$ acts as a
lower bound on $\log p(\x;\Vtheta)$ for all values of $\Vtheta$. Maximising over $\Vtheta$
yields $\prm{\Vtheta^{(k+1)}}$, which achieves a higher ELBO value than
$\prm{\Vtheta^{(k)}}$ and corresponds to a new, generally higher log-likelihood (purple
curve).}\label{f:elbo2}
\end{figure}

To summarise, each iteration of this general procedure consists of two steps. In the first
step, one finds the optimal auxiliary distribution given the current parameters---that is, the
distribution that maximises the ELBO at the current $\Vtheta$:
\begin{myequation}
  \func{q^{(k)}}=\argmax{\vr{q}}{\mL(\vr{q},\ox, \prm{\Vtheta^{(k)}})}\,=\,p(\z|\ox;\prm{\Vtheta^{(k)}}).
\end{myequation}
In the second step, one updates the parameters by maximising the ELBO with $q$ now
fixed at $\func{q^{(k)}}$, which reduces to maximising the expected complete-data
log-likelihood:
\begin{myequation}
  \prm{\Vtheta^{(k+1)}}=\argmax{\vrt}{\mL(\func{q^{(k)}},\ox,\vrt)}=\argmax{\vrt}{Q(\func{q^{(k)}},\ox,\vrt)}.
\end{myequation}
This two-step structure---alternating between updating $q$ to tighten the bound and
updating $\Vtheta$ to maximise it---is the conceptual skeleton of the EM algorithm. The
first step corresponds to the E-step (Expectation), in which the posterior over latent
variables is computed; the second corresponds to the M-step (Maximisation), in which the
parameters are updated. We will make this connection precise in the next section, where
we also address the question of whether and when this procedure guarantees monotone
improvement of the observed-data log-likelihood.

\subsection*{Monotonicity and Its Limits}

Having described the two-step iterative procedure, it is natural to ask whether it guarantees
monotone improvement of the log-likelihood at each iteration---that is, whether
$\log p(\ox;\prm{\Vtheta^{(k+1)}}) \geq \log p(\ox;\prm{\Vtheta^{(k)}})$ is guaranteed to
hold. The answer, in the general formulation presented so far, is no, and understanding
precisely why illuminates the additional ingredient that the EM algorithm contributes.

By construction, the following chain of inequalities holds at each iteration:
\begin{myalign}
\log p(\ox;\prm{\Vtheta^{(k)}})&\geq \mL(\func{q^{(k)}},\ox,\prm{\Vtheta^{(k)}}),\\
\log p(\ox;\prm{\Vtheta^{(k+1)}})&\geq \mL(\func{q^{(k)}},\ox,\prm{\Vtheta^{(k+1)}})\geq \mL(\func{q^{(k)}},\ox,\prm{\Vtheta^{(k)}}).
\end{myalign}
The first line states that the ELBO is a lower bound on the log-likelihood at the current
parameters---this is simply the fundamental bound established earlier. The second line
combines two facts: the ELBO is a lower bound on the log-likelihood at any parameter
value (first inequality), and $\prm{\Vtheta^{(k+1)}}$ was chosen to maximise the ELBO
with $\func{q^{(k)}}$ fixed, so it must achieve at least as high an ELBO value as
$\prm{\Vtheta^{(k)}}$ (second inequality). Together, these inequalities confirm that the
ELBO value at $\ox$ has increased from one iteration to the next. However, they do
\emph{not} in general guarantee that the true log-likelihood has increased, i.e., that
\begin{myequation}
\log p(\ox;\prm{\Vtheta^{(k+1)}})\geq \log p(\ox;\prm{\Vtheta^{(k)}}).
\end{myequation}

The reason is subtle but important. The ELBO at $\prm{\Vtheta^{(k+1)}}$ is only a lower
bound on $\log p(\ox;\prm{\Vtheta^{(k+1)}})$, not an equality. Even though the ELBO has
increased, the gap $\mK(\func{q^{(k)}},\ox,\prm{\Vtheta^{(k+1)}}) =
\dkl{\func{q^{(k)}}(\z)}{p(\z|\ox;\prm{\Vtheta^{(k+1)}})}$ at the new parameters may be
larger than the gap at the old parameters. If this additional gap is large enough, it could
more than offset the improvement in the ELBO, leaving the true log-likelihood lower than
before. This pathological situation is illustrated in Figure~\ref{f:elbo3}: despite the ELBO
increasing from $\mL(\func{q^{(k)}},\ox,\prm{\Vtheta^{(k)}})$ to
$\mL(\func{q^{(k)}},\ox,\prm{\Vtheta^{(k+1)}})$, the log-likelihood at the new parameters
$\log p(\ox;\prm{\Vtheta^{(k+1)}})$ is lower than at the old ones.

\begin{figure}[h]
\begin{center}
\begin{tikzpicture}[scale=1]
\begin{axis}[
    clip=false,axis lines=middle,
    xmin=-2, xmax=16, xlabel = $\x$,
    ymin=-3, ymax=1,  ylabel={},
    xticklabels = {},
    yticklabels = {},
    extra x ticks={4},
    extra x tick style={grid=major,
        tick label style={anchor=north west}},
    extra x tick labels={$\ox$},
    legend style={
        legend pos=outer north east,
    },]
\addplot [xkcdCadetBlue, line width=.4mm, smooth, tension=.9] coordinates {(-2, -2.9) (0,-2.7) (2,-1.6) (4,-.4) (6,-1.3) (8,-2) (10,-.7) (12,-1.5) (14,-2.3)};
\addlegendentry{\footnotesize  $\log p(\x;\prm{\Vtheta^{(k)}})$}
\addplot [xkcdRed, smooth, tension=.9] coordinates {(-2, -3) (0,-3) (2,-2) (4,-1.5) (6,-1.8) (8,-2.5) (10,-2) (12,-2.2) (14,-2.7)};
\addlegendentry{\footnotesize  $\mL(\func{q^{(k)}},\x,\prm{\Vtheta^{(k)}})$}
\addplot [xkcdBrownyOrange,smooth, tension=.9] coordinates {(-2, -2.4) (0,-2.7) (2,-2.6) (4,-2.7) (6,-2.1) (8,-1.8) (10,-1.4) (12,-2.1) (14,-2.3)};
\addlegendentry{\footnotesize  $\mL(\func{q^{(k)}},\x, \prm{\Vtheta'})$}
\addplot [xkcdFlatGreen,smooth, tension=.9] coordinates {(-2,-2.7) (0,-2.4) (2,-1.7) (4,-.8) (6,-.6) (8,-1.5) (10,-2) (12,-1.7) (14,-2.5)};
\addlegendentry{\footnotesize  $\mL(\func{q^{(k)}},\x, \prm{\Vtheta^{(k+1)}})$}
\filldraw (4,-1.5) [xkcdRed]circle (2pt)node{};
\filldraw (-.5,-1.1) node[left, font=\footnotesize] {$\mL(\func{q^{(k)}},\ox,\prm{\Vtheta^{(k)}})$};
\draw[-stealth]  (-.5,-1.1) -- (4,-1.5);
\filldraw (4,-.8) [xkcdFlatGreen]circle (2pt)node{};
\filldraw (7.5,-.5) node[right, font=\footnotesize] {$\mL(\func{q^{(k)}},\ox,\prm{\Vtheta^{(k+1)}})$};
\draw[-stealth]  (7.5,-.5) -- (4,-.8);
\filldraw (4,-2.7) [xkcdBrownyOrange]circle (2pt)node{};
\filldraw (5,-3.2) node[right, font=\footnotesize] {$\mL(\func{q^{(k)}},\ox,\prm{\Vtheta'})$};
\draw[-stealth]  (5,-3.2) -- (4,-2.7);
\addplot [xkcdBarneyPurple, line width=.4mm, smooth, tension=.9] coordinates {(-2, -2.5) (0,-2.2) (2,-1.2) (4,-.7) (6,-.3) (8,-1.3) (10,-1.5) (12,-1.6) (14,-2.2)};
\addlegendentry{\footnotesize  $\log p(\x;\prm{\Vtheta^{(k+1)}})$}
\end{axis}
\end{tikzpicture}
\end{center}
\caption{An example of an iteration in which the ELBO increases but the log-likelihood does
not: the gap $\mK$ at the new parameters is sufficiently large to offset the improvement in
the lower bound.}\label{f:elbo3}
\end{figure}

Monotone increase of the log-likelihood can be guaranteed if, in the first step, the
auxiliary distribution is set to the exact posterior $q(\z) = p(\z|\ox;\prm{\Vtheta^{(k)}})$.
In this case, the gap $\mK$ at the current parameters is exactly zero, meaning the ELBO
touches the log-likelihood from below at $\prm{\Vtheta^{(k)}}$. Any subsequent increase
in the ELBO through the M-step then translates directly into an increase in the
log-likelihood, since the starting point of the comparison is exact---the ELBO and the
log-likelihood coincide at $\prm{\Vtheta^{(k)}}$, so any improvement in the ELBO at the
new parameters implies a strictly higher lower bound, which in turn lies below a
log-likelihood that can only be at least as large. This is precisely the mechanism that the
EM algorithm exploits, and we will formalise it in the next section.

\subsection*{Extension to a Full Dataset and Amortisation}

The single-observation analysis generalises straightforwardly to a full dataset $\X =
\{\vlxu,\ldots,\vlxn\}$ of $n$ observations. Since the observations are independent, the
dataset log-likelihood decomposes as a sum, and the same decomposition into ELBO and
KL gap applies term-wise. Introducing a separate auxiliary distribution $q_i(\z)$ for each
observation $\vlxi$, we obtain:
\begin{myalign}
\log p(\X ;\Vtheta)&=\sum_{i=1}^n \log p(\vlxi;\Vtheta)\\
&=\sum_{i=1}^n\ev{q_i(\z)}{\log\frac{p(\vlxi,\dz;\Vtheta)}{p(\dz|\vlxi;\Vtheta)}}\\
&=\sum_{i=1}^n\mL(q_i,\vlxi, \Vtheta)+\sum_{i=1}^n\mK(q_i,\vlxi, \Vtheta)\\
&\defeq\mL(\mathbf{q},\vlX,\Vtheta)+\mK(\mathbf{q},\X, \Vtheta),
\end{myalign}
where $\mathbf{q}=(q_1,\ldots,q_n)$ is a collection of auxiliary distributions, one for each
data point. The total ELBO on the dataset log-likelihood is therefore the sum of the
individual ELBOs. Equality between the total ELBO and the dataset log-likelihood holds if
and only if each $q_i(\z) = p(\z|\vlxi;\Vtheta)$, i.e., when every auxiliary distribution
coincides with the corresponding exact posterior.

The iterative procedure extends accordingly. In the first step, one finds the optimal
distribution $\func{q_i^{(k)}}(\z)$ separately for each observation $\vlxi$ by minimising
$\mK(\func{q_i^{(k)}},\vlxi,\prm{\Vtheta^{(k)}})$, which again yields the exact posterior
$p(\z|\vlxi;\prm{\Vtheta^{(k)}})$ as the optimal choice. In the second step, one maximises
the total ELBO over $\Vtheta$:
\begin{myequation}
    \mL(\func{\mathbf{q}^{(k)}}, \X,\Vtheta)=\sum_{i=1}^n\mL(\func{q_i^{(k)}},\vlxi,\Vtheta),
\end{myequation}
which, as in the single-observation case, reduces to maximising the total expected
complete-data log-likelihood under the collection $\func{\mathbf{q}^{(k)}}$.

This multi-observation formulation, while conceptually clean, raises a practical concern of
considerable importance. Maintaining and optimising a separate distribution $q_i(\z)$ for
each of the $n$ observations requires memory and computation that scale linearly with the
dataset size. In large-scale settings---where $n$ may be in the millions---this is
prohibitive. A practically important and widely adopted simplification is to replace the
entire collection $\{q_1,\ldots,q_n\}$ with a single shared conditional distribution
$q(\z|\x)$, setting $q_i(\z) = q(\z|\vlxi)$ for all $i$. Rather than storing and optimising a
separate distribution for each data point, one learns a single parametric mapping from
observations to distributions over latent variables. This strategy is known as
\alert{amortisation}: the inference cost is amortised across all data points by sharing a
common functional form for $q(\z|\x)$, whose parameters are optimised jointly with those
of the generative model. The quality of the approximation is then limited by the
expressiveness of the chosen parametric family for $q(\z|\x)$, rather than by
computational constraints. This trade-off---sacrificing exact posterior inference for
scalability---is a defining feature of modern deep generative models. Amortised inference
of this kind is the foundation of Variational Autoencoders, where $q(\z|\x)$ is
parameterised by a neural network---the encoder---trained jointly with the decoder network
that parameterises $p(\x|\z)$. The encoder learns to map any observation directly to the
parameters of an approximate posterior, enabling both efficient training and fast inference
at test time.

%\end{document}